\documentclass[11pt]{article}

\usepackage[T1]{fontenc}
\usepackage{amsmath}
\usepackage{amssymb}
\usepackage{booktabs}
\usepackage[hidelinks]{hyperref}

\newcommand{\claimstatus}[2]{\textbf{#1 [#2]}}

\title{Order-Separating Exponent Families for Classical Integer
Factorization: Algorithms, Barriers, and Reproducible Evaluation}
\author{MOSEF Research Program}
\date{\today}

\begin{document}
\maketitle

\begin{abstract}
This initial manuscript defines an evidence-driven program for evaluating
order-separating exponent families in classical integer factorization. It
introduces no new factoring theorem. In particular, the existence of a
universally constructible polynomial order-separating family remains open.
The present contribution is a reproducibility and claim-tracking foundation
for later algorithms, counterexamples, restricted theorems, and barriers.
\end{abstract}

\section{Introduction}

Let \(N\) be a positive integer represented in binary. The project asks whether
order information from one or more algebraic groups can expose a nontrivial
factor without circular access to the factorization. General classical
polynomial-time factorization is a research target, not a premise.

The immediate objective is methodological: every mathematical statement is
linked to a status and evidence record, finite experiments are not promoted to
universal theorems, and complexity is measured in binary input length.

\section{Model and definitions}

\claimstatus{DEF-001}{DEFINITION}. The input length is
\[
  m=\left\lceil\log_2 N\right\rceil.
\]
Thus a running time polynomial in \(N\) is not, merely for that reason,
polynomial in the input length.

\claimstatus{DEF-002}{DEFINITION}. Let \(P(N)\) be the set of distinct prime
divisors of \(N\). For \(\gcd(g,N)=1\) and \(d>0\), define
\[
  D_N(g,d)=\{p\in P(N):\operatorname{ord}_p(g)\mid d\}.
\]
The pair \((g,d)\) is an order separator when \(D_N(g,d)\) is nonempty and is
not all of \(P(N)\).

A polynomial order-separating family must bound not only the number of bases
and exponents, but also construction cost, exponent bit lengths, evaluation
cost, and recursive factor extraction in \(\operatorname{poly}(m)\) bit
operations.

\section{Inspected baseline evidence}

\claimstatus{EXT-001}{PROVED, IMPORTED}. Agrawal, Kayal, and Saxena give an
unconditional deterministic polynomial-time primality test
\cite{agrawal2004primes}. This decision result must not be restated as a
factor-finding algorithm.

\claimstatus{EXT-002}{PROVED, IMPORTED}. Lenstra and Pomerance prove a
probabilistic complete-factorization algorithm with expected running time
\(L_N[1/2,1+o(1)]\) \cite{lenstra1992rigorous}. This is a rigorous
subexponential bound, not a polynomial bound in \(m\).

Lenstra's elliptic-curve method motivates varying the group channel, while the
expected-time expression highlighted in the inspected abstract is explicitly
conjectural \cite{lenstra1987elliptic}. It therefore supplies motivation, not a
separation theorem for MOSEF.

\section{Open research target}

\claimstatus{OPEN-001}{OPEN}.

No proof in this repository establishes a factorization-independent
construction for every composite \(N\). In
particular, no construction is known here that satisfies every family-size,
exponent-length, evaluation, separation, and polynomial-time requirement.

The research program will first build trusted reference algorithms and exact
test vectors, then formalize all MOSEF failure branches and search for small
counterexamples before attempting a restricted theorem.

\section{Baseline algorithm suite}

The M1 suite contains exact modular exponentiation, trial division,
perfect-power detection, exact trial-division primality for validation,
bounded Pollard rho, Pollard \(p-1\) stage 1, a scoped Williams-style
\(p+1\) stage 1 using the Lucas sequence \(V_k(P,1)\), and per-item batch
GCD semantics.

Python supplies the arbitrary-precision reference behavior. Rust supplies the
authoritative implementation only on its documented \texttt{u64} domain, with
\texttt{u128} products preventing modular-multiplication overflow. An
independent C\# \texttt{BigInteger} program verifies selected operations. The
current batch-GCD function is a semantic baseline and is not a product-tree
performance claim.

\claimstatus{EMP-001}{EMPIRICAL}. On the finite canonical corpus in
\path{schemas/baseline-vectors-v1.json}, 58 cross-language comparisons
passed. Rust agreed with Python for every registered operation; C\# also agreed
for modular exponentiation, validation primality, trial factors, and per-item
batch GCDs. This finite agreement neither proves asymptotic complexity nor
establishes correctness outside the specified domains.

\section{Reproducibility}

The result envelope is versioned by a JSON schema stored in the repository's
\texttt{schemas} directory. Each record contains the repository commit, UTC
timestamp, host and toolchain metadata, deterministic seed, input and output
hashes, timing fields, memory metadata, and explicit status. Raw records are
immutable; corrections must supersede rather than edit them.

The dependency-free foundation gates are:
\begin{verbatim}
python scripts/validate_foundation.py
python -m unittest discover -s tests -v
cargo fmt --all --check
cargo clippy --workspace --all-targets --all-features -- `
  -D warnings
cargo test --workspace --all-features
dotnet build verification/csharp/MosefVerifier.csproj
python scripts/check_baseline_differential.py
\end{verbatim}

\section{Limitations}

No POSF construction, restricted-class theorem, counterexample corpus, or
performance experiment is claimed here. Trial-division primality is intentionally
not polynomial in binary input length, the Rust algorithms are bounded to
\texttt{u64}, and finite differential tests are not proofs for untested inputs.
The baseline literature set is not a current-state-of-the-art survey. The
current four-page manuscript is an initialization artifact, not a
publication-ready contribution; it has passed XeLaTeX compilation, citation
resolution, warning scanning, and page-by-page visual inspection.

\section{Conclusion}

The M0 foundation makes claim status, source inspection, and experiment
metadata explicit, and M1 provides a differentially checked small-input
baseline. The next milestone formalizes all MOSEF/POSF branches and the basic
separator lemma; universal POSF existence remains open.

\bibliographystyle{plain}
\bibliography{paper/references}

\end{document}
